\documentclass[11pt]{article}

\usepackage[T1]{fontenc}
\usepackage{lmodern}
\usepackage{amsmath,amssymb,amsthm}
\usepackage[margin=1in]{geometry}
\usepackage{microtype}
\usepackage[hidelinks]{hyperref}
\hypersetup{
  pdftitle={Almost everywhere divergence of polynomial interpolation with sublinear excess degree}
}

\newtheorem{theorem}{Theorem}
\newtheorem{lemma}[theorem]{Lemma}
\newtheorem{proposition}[theorem]{Proposition}
\theoremstyle{definition}
\newtheorem{definition}[theorem]{Definition}
\numberwithin{equation}{section}
\DeclareMathOperator{\supp}{supp}
\DeclareMathOperator{\sgn}{sgn}
\DeclareMathOperator{\arctanh}{arctanh}
\newcommand{\R}{\mathbb R}
\newcommand{\C}{\mathbb C}
\newcommand{\Z}{\mathbb Z}
\newcommand{\one}{\mathbf 1}
\newcommand{\norm}[1]{\left\lVert #1\right\rVert}

\title{Almost everywhere divergence of polynomial interpolation\\
with sublinear excess degree}
\author{Qiyuan Gu\\
University of Chicago\\
\href{mailto:phoenix1203@uchicago.edu}{\texttt{phoenix1203@uchicago.edu}}}
\date{}

\begin{document}
\maketitle

\begin{abstract}
Let \(X_n\) be an arbitrary set of \(n\) distinct points in \([-1,1]\),
and let \(r_n\) be nonnegative integers with \(r_n=o(n)\).
We prove that there is a real continuous function \(f\) such that every
sequence of polynomials \(p_n\), of degree at most \(n+r_n\), satisfying
\(p_n=f\) on \(X_n\), has
\(\limsup_{n\to\infty}|p_n(x)|=\infty\) for almost every \(x\in[-1,1]\).
The exceptional null set may depend on the sequence \((p_n)\).
This extends the Erd\H{o}s--V\'ertesi theorem on almost-everywhere
unboundedness of Lagrange interpolation for arbitrary node arrays to
nonunique interpolation with \(o(n)\) excess degree, and in particular
answers Erd\H{o}s Problem 1152 in a stronger form. The proof combines
localized functions in Bernstein spaces with uniform asymptotics for
weighted Christoffel--Darboux kernels. A finite interpolation construction
and the Baire category theorem give a single function valid for every
admissible interpolation sequence.
\end{abstract}

\section{Introduction}

Write \(\Pi_d\) for the space of real polynomials of degree at most \(d\).
All measures of subsets of \([-1,1]\) below are Lebesgue measures.
The uniform norm on this interval is denoted by \(\norm{\cdot}_\infty\).

\begin{theorem}\label{thm:main}
For each \(n\ge1\), let \(X_n\subset[-1,1]\) consist of \(n\) distinct
points. Suppose that \(r_n\ge0\) are integers and \(r_n/n\to0\).
There exists \(f\in C([-1,1];\R)\) such that every sequence satisfying
\[
 p_n\in\Pi_{n+r_n},
 \qquad p_n(x)=f(x)\quad(x\in X_n)
\]
also satisfies
\begin{equation}\label{eq:divergence}
 \limsup_{n\to\infty}|p_n(x)|=\infty
 \quad\text{for almost every }x\in[-1,1].
\end{equation}
\end{theorem}

Erd\H{o}s and V\'ertesi proved that for every triangular array of
distinct nodes in \([-1,1]\), there is a continuous function whose
Lagrange interpolation polynomials are unbounded almost everywhere
\cite{EV80,EV81}. Theorem~\ref{thm:main} extends this arbitrary-node
result to nonunique interpolation: degrees up to \(n+r_n\), with
\(r_n=o(n)\), are allowed, and the same continuous function works for
every admissible sequence of interpolants. This also answers the question
in \cite[Problem~2.42]{Favorite99}, recorded as Erd\H{o}s Problem 1152
\cite{Erdos1152}, in the stronger form \eqref{eq:divergence}.
Indeed, for the degree condition \(\deg p_n<(1+\varepsilon_n)n\),
set \(r_n=\max\{0,\lceil\varepsilon_n n\rceil-1\}\); then
\(r_n=o(n)\) and \(\deg p_n\le n+r_n\).

The contrast with fixed positive relative excess is provided by Erd\H{o}s,
Kro\'o, and Szabados \cite{EKS89}: for suitable node arrays and each fixed
\(\varepsilon>0\), every continuous function admits uniformly convergent
interpolants of degree at most \(\lfloor(1+\varepsilon)n\rfloor\).
For background on divergence phenomena in Lagrange interpolation, see
Szabados and V\'ertesi \cite{SV90}.

The two analytic ingredients are the construction in
Olevskii and Ulanovskii \cite[Lemma 4.1]{OU} and the uniform kernel
asymptotics of Kriecherbauer, Schubert, Sch\"uler, and Venker
\cite[Theorem 1.3 and Theorem 1.8 (bulk universality)]{KSSV}.
Section~\ref{sec:kernel} records the kernel estimates in the normalization
used below, while Section~\ref{sec:arbitrary} derives the corresponding
local external field from a weak limit of the node measures.
We first obtain interpolation data that force large values on a fixed
fraction of a local interval, uniformly over all admissible polynomial
corrections. We then combine finitely many such data sets to rule out
common sets of bounded values. The final step is a category argument
in \(C([-1,1];\R)\).

\section{Localized Bernstein functions and polynomial corrections}
\label{sec:bernstein}

For \(\sigma>0\), let \(B_\sigma\) denote the space of entire functions
of exponential type at most \(\sigma\) that are bounded on \(\R\).
Norms of functions in \(B_\sigma\) are taken on \(\R\).

\begin{lemma}\label{lem:bernstein}
There are absolute constants \(c,C>0\) with the following property.
For every sufficiently large integer \(\rho\), put
\[
 R_0=\rho-\sqrt{\rho},
 \qquad \gamma=R_0^{-1/9}.
\]
For every finite set \(\Gamma\subset[-\rho,\rho)\) with
\(\#\Gamma\le2\rho+1\), there is an entire function \(g\), real on
\(\R\), such that
\begin{align}
 &g\in B_{\pi-4\gamma},
 \qquad \max_{\Gamma}|g|\le1,
 \qquad \norm{g}_\infty\le C\log\rho, \label{eq:bernstein-norm}\\
 &g(u_0)\ge c\log\rho
 \quad\text{for some }|u_0|\le R_0, \label{eq:bernstein-peak}\\
 &|g(u)|\le C(|u|-R_0)^{-3}
 \quad (|u|\ge\rho). \label{eq:bernstein-tail}
\end{align}
For each fixed \(\rho\), these choices \(g=g_\Gamma\) satisfy
\[
 \sup_\Gamma\norm{g_\Gamma}_{L^1(\R)}\le C_\rho,
 \qquad
 \sup_\Gamma\int_{|u|>T}|g_\Gamma(u)|\,du
 \le\frac{C}{(T-R_0)^2}\quad(T\ge\rho),
\]
where the suprema range over the finite sets \(\Gamma\) specified above.
There is also an interval, of length bounded below by an absolute
positive constant, contained in \((-\rho,\rho)\setminus\Gamma\), on
which \(g\ge c'\log\rho\), with an absolute constant \(c'>0\).
\end{lemma}

\begin{proof}
Adapt the construction of \cite[Lemma~4.1, pp.~184--185]{OU} by reducing
its initial bandwidth from \(\pi-4R_0^{-1/9}\) to
\(\pi-8R_0^{-1/9}\). The fourth-power sinc factor below then leaves
the required margin \(4\gamma\) below \(\pi\).
Put \(w=\pi-8\gamma\) and
\[
 \Lambda=\Gamma+2R_0\Z.
\]
Since \(\Gamma\) is finite, \(\Lambda\) is uniformly discrete. Being
periodic, its lower and upper uniform densities coincide; write their
common value as \(D\). Then
\[
 D\le\frac{2\rho+1}{2R_0}\le1+CR_0^{-1/2}.
\]
Write \(K_s(\Lambda,B_w)\) for the sampling constant. If \(\Lambda\)
is a sampling set for \(B_w\), then \cite[Theorem~C(i), p.~179]{OU} gives
\(D>w/\pi\). Hence \(D\Lambda\) has density one and
\cite[Theorem~1(i), p.~180]{OU} gives
\[
 K_s(\Lambda,B_w)
 =K_s(D\Lambda,B_{w/D})
 \ge C\log\frac{\pi}{\pi-w/D}
 \ge c\log R_0,
\]
since, for large \(\rho\),
\[
 0<\pi-\frac wD=\frac{\pi D-w}{D}
 \le C\bigl(R_0^{-1/2}+\gamma\bigr)\le C\gamma.
\]
If \(\Lambda\) is not a sampling set,
its sampling constant is infinite, so the same lower bound is immediate.
Thus one can choose \(g_1\in B_w\), bounded by one on \(\Lambda\),
whose norm is a fixed constant times \(\log R_0\).
Choose a point at which its absolute value is at least half this norm.
A translation by a multiple of \(2R_0\) places this point at
\(u_0\in[-R_0,R_0]\), without changing the bound on \(\Lambda\).
Multiply by a constant of modulus one so that \(g_1(u_0)>0\), then
replace \(g_1(z)\) by
\((g_1(z)+\overline{g_1(\bar z)})/2\). This entire function is real on
\(\R\), retains the value at \(u_0\), and satisfies the same bandwidth
and upper bounds.

Multiply this function by
\[
 \left(\frac{\sin(\gamma(u-u_0))}
 {\gamma(u-u_0)}\right)^4.
\]
The resulting function has bandwidth at most \(\pi-4\gamma\) and
retains the node bound and the value at \(u_0\).
For \(|u|\ge\rho\), the bound
\[
 C\log R_0\,\gamma^{-4}(|u|-R_0)^{-4}
 \le C(|u|-R_0)^{-3}
\]
follows from
\(\log R_0\,R_0^{4/9}/\sqrt{\rho}=O(1)\).
This proves \eqref{eq:bernstein-norm}--\eqref{eq:bernstein-tail}
and the assertions about integrability.
Finally, Bernstein's inequality gives
\(\norm{g'}_\infty\le\pi C\log\rho\).
Thus the peak contains an interval of fixed positive length on which
\(g\ge c'\log\rho>1\). This interval contains no point of \(\Gamma\)
and lies inside the original box.
\end{proof}

For a finite set \(Y\subset\R\), define
\[
 P_Y(x)=\prod_{y\in Y}(x-y).
\]

\begin{lemma}[Alternating peaks]\label{lem:signs}
Let \(Y\) consist of \(N\ge1\) distinct real points, let \(H>0\),
and let \(d\ge0\) be an integer. Suppose that disjoint intervals
\(J_1,\ldots,J_K\), in increasing order and missing \(Y\), and
polynomials \(B_1,\ldots,B_K\in\Pi_{N-1}\) satisfy
\[
 \max_{y\in Y}\sum_j|B_j(y)|\le2,
 \qquad
 B_j(x)>2H+\sum_{i\ne j}|B_i(x)|\quad(x\in J_j).
\]
There are data \(v_y\in[-1,1]\) on \(Y\) such that every
\(p\in\Pi_{N+d}\) interpolating these data satisfies \(|p|>H\)
throughout at least \((K-d-1)/2\) of the intervals \(J_j\).
\end{lemma}

\begin{proof}
Choose signs \(\varepsilon_j\) so that
\(\sgn(\varepsilon_j/P_Y)|_{J_j}=(-1)^j\), and put
\(b=\frac12\sum_j\varepsilon_jB_j\).
Then \(|b|\le1\) on \(Y\), while peak dominance gives
\begin{equation}\label{eq:alternating-data}
 |b|>H,
 \qquad
 \sgn(b/P_Y)|_{J_j}=(-1)^j.
\end{equation}
Set \(v_y=b(y)\). Any interpolant in \(\Pi_{N+d}\) has the form
\(p=b+P_Yq\) with \(\deg q\le d\).
Call \(J_j\) a low interval if it contains a point \(x_j\) with
\(|p(x_j)|\le H\). At such a point the sign of \(q(x_j)\) is opposite
to that of \(b(x_j)/P_Y(x_j)\).
If there are \(B\) low intervals, the resulting signs alternate at
least \(2B-K-1\) times. Indeed, deleting \(K-B\) terms from a fully
alternating sequence removes at most \(2(K-B)\) sign changes.
Consequently \(2B-K-1\le d\), so
\begin{equation}\label{eq:low-interval-count}
 B\le\frac{K+d+1}{2}.
\end{equation}
The asserted bound follows.
\end{proof}

\section{An explicit external field and its polynomial kernel}
\label{sec:kernel}

For a finite positive measure \(\mu\), use the logarithmic-potential
convention
\[
 V_\mu(x)=\int\log|x-y|\,d\mu(y).
\]
Define
\begin{equation}\label{eq:field}
 Q(u)=\frac12\bigl((1+u)\log(1+u)+(1-u)\log(1-u)\bigr),
 \qquad |u|<1.
\end{equation}
Then \(Q(0)=0\), \(Q'(u)=\arctanh u\), and
\(Q''(u)=(1-u^2)^{-1}\).

\begin{lemma}\label{lem:equilibrium}
For \(0<t<1\), set \(a_t=\sqrt{1-(1-t)^2}\).
The equilibrium measure of mass \(t\) in the external field \(Q\)
on \([-1,1]\) has density
\begin{equation}\label{eq:density}
 \rho_t(u)=\frac1\pi
 \arccos\frac{1-t}{\sqrt{1-u^2}}\,
 \one_{\{|u|<a_t\}}.
\end{equation}
Thus \(Q-V_{\rho_t}\) is constant on \([-a_t,a_t]\) and is no smaller
elsewhere in \([-1,1]\).
\end{lemma}

\begin{proof}
Differentiation with respect to \(t\), away from \(|u|=a_t\), gives
\begin{equation}\label{eq:density-derivative}
 \partial_t\rho_t(u)=
 \frac{\one_{\{|u|<a_t\}}}{\pi\sqrt{a_t^2-u^2}}.
\end{equation}
This is the arcsine probability density on \([-a_t,a_t]\), so
\(\rho_t\) has total mass \(t\).
Its Cauchy transform can be computed by integrating the transforms
of these arcsine measures. For \(0<x<a_t\), its real boundary value is
\begin{equation}\label{eq:cauchy}
 \int_0^{1-\sqrt{1-x^2}}\frac{ds}{\sqrt{x^2-a_s^2}}
 =\arctanh x=Q'(x).
\end{equation}
One may justify the calculation first in the upper half-plane by
Fubini's theorem and then take boundary values.
For \(a_t<x<1\), the upper limit in \eqref{eq:cauchy} is \(t\),
and the integral is strictly smaller than \(Q'(x)\).
Symmetry gives the corresponding statements on the negative
half-axis. These identities prove the equilibrium equalities and
inequalities.
\end{proof}

Put \(A=[-1/4,1/4]\). We will use the following consequences of
\cite{KSSV}.

\begin{lemma}[Uniform weighted kernel estimates]\label{lem:kssv-kernel}
Fix \(0<t<1\) sufficiently close to one that \(A\Subset(-a_t,a_t)\),
and choose \(a_t<b<1\). There are intervals
\[
 A\Subset J_0\Subset J_1\Subset(-b,b)
\]
and a bounded complex neighborhood \(\Omega\) of \([-b,b]\), with
\(\overline\Omega\subset\{z:|\operatorname{Re}z|<1\}\), and
\(\delta>0\) with the following property.
Let \(\mathcal U\) consist of the bounded holomorphic functions \(F\)
on \(\Omega\), real on \(\Omega\cap\R\), for which
\(\norm{F-Q}_\Omega<\delta\), where
\(\norm{F}_\Omega=\sup_{z\in\Omega}|F(z)|\).
Let \(s\to\infty\), \(m=\lfloor ts\rfloor\), and let \(Q_s\) belong
to \(\mathcal U\). If \(p_0,\ldots,p_{m-1}\) are orthonormal for
\(e^{-2sQ_s(x)}\,dx\) on \([-b,b]\), set
\begin{equation}\label{eq:kernel}
 K_s(x,z)=e^{-sQ_s(x)-sQ_s(z)}
 \sum_{j=0}^{m-1}p_j(x)p_j(z),
 \qquad c_s=\frac{s}{2}.
\end{equation}
Let \(\beta_{s,y}\) be twice the density at \(y\) of the equilibrium
measure of mass \(m/s\) for \(Q_s\). Then, uniformly for the indicated
family of external fields,
\begin{align}
 \frac{K_s(y+u/c_s,y+v/c_s)}{c_s}
 &=\frac{\sin\bigl(\pi\beta_{s,y}(u-v)\bigr)}{\pi(u-v)}+o(1),
 \label{eq:sine-limit}\\
 \frac{|K_s(x,z)|}{c_s}
 &\le \frac{C}{1+s|x-z|}
 \qquad (x,z\in J_1), \label{eq:bulk-kernel}\\
 \frac{|K_s(x,z)|}{c_s}
 &\le Cs^{-5/6}
 \qquad\bigl(x\in[-b,b]\setminus J_0,\
 z\in A+O(s^{-1})\bigr). \label{eq:edge-kernel}
\end{align}
In \eqref{eq:sine-limit}, \(y\in A\), \(u,v\) range over bounded
sets, and the quotient at \(u=v\) is interpreted by continuity.
\end{lemma}

\begin{proof}
In the notation of \cite{KSSV}, take
\[
 N=m,\qquad J=[-b,b],\qquad V=V_s:=\frac{2s}{m}Q_s.
\]
Then \(e^{-mV_s}=e^{-2sQ_s}\), and their kernel \(K_{N,V}\) is
exactly \(K_s\). The base field \(V_0=2Q/t\) satisfies their
assumptions (GA)\(_2\): it is analytic near the compact interval
\(J\), has \(V_0''\ge2/t>0\), and its equilibrium probability
measure is \(\rho_t/t\), supported on \([-a_t,a_t]\Subset J\).

Choose \(\Omega\) as in \cite[Remark~1.1]{KSSV} and a fixed
sup-norm neighborhood of \(V_0\) supplied by their Theorem~1.3(b).
Since \(m/s\to t\), decreasing \(\delta\) and then taking \(s\)
large ensures that every \(V_s\) lies in this same neighborhood.
By their Lemma~2.4 the support endpoints depend continuously on
\(V\). Thus \(J_0,J_1\) can be chosen inside every support, with
fixed positive distances between \(A\), their boundaries, the
support endpoints, and \(\pm b\). Cauchy's estimates also preserve
a positive lower bound for \(V_s''\). In particular, the density
factor \(G_V\) in their (1.14)--(1.15) stays positive and bounded
on the corresponding rescaled intervals.

Let \(\psi_s\) be the equilibrium probability density for \(V_s\).
Their affine map \(\lambda_{V_s}\) takes \([-1,1]\) onto the
equilibrium support, and their rescaled density satisfies
\[
 \rho_{V_s}(\xi)=\lambda_{V_s}'\,
                  \psi_s(\lambda_{V_s}(\xi)).
\]
Then \(\rho_s=(m/s)\psi_s\) is the equilibrium density of mass \(m/s\)
for \(Q_s\), and
\[
 m\psi_s(y)=s\rho_s(y)=c_s\beta_{s,y}.
\]
In \cite[Theorem~1.8]{KSSV}, take the base point
\(\lambda_{V_s}^{-1}(y)\) and microscopic variables
\(\beta_{s,y}u\), \(\beta_{s,y}v\).
The preceding identities give \eqref{eq:sine-limit}, with the
uniform error in part (b) of that theorem.

For \eqref{eq:bulk-kernel}, the bulk vector \(k\) in
\cite[Theorem~1.3(a), (1.25)]{KSSV} is uniformly bounded.
The difference quotient therefore gives \(O(|x-z|^{-1})\) when
\(|x-z|\ge s^{-1}\). For \(|x-z|<s^{-1}\), the diagonal bound
\(K_s(x,x)=O(s)\) from Theorem~1.8 and Cauchy--Schwarz give the
required \(O(s)\) bound.
For \eqref{eq:edge-kernel}, the endpoint components of the vector
\(k\) in that theorem have the forms, up to bounded multipliers,
\[
 m^{1/6}\operatorname{Ai}(\zeta),\qquad
 m^{-1/6}\operatorname{Ai}'(\zeta),\qquad |\zeta|=O(m^{2/3}).
\]
The real Airy bounds \(|\operatorname{Ai}(\zeta)|\le C\) and
\(|\operatorname{Ai}'(\zeta)|\le C(1+|\zeta|)^{1/4}\) give
\(\norm{k}=O(m^{1/6})\) there. The bulk and exterior formulas give
bounded vectors away from the endpoints. At \(z\in A+O(s^{-1})\)
the vector is bounded, while \(|x-z|\) is bounded below.
Thus (1.25), including its \(O(m^{-1})\) remainder, gives
\(|K_s(x,z)|=O(m^{1/6})\). Since \(m\asymp s\) and \(c_s=s/2\),
this proves \eqref{eq:edge-kernel}.
The neighborhood and error bounds are uniform by Theorem~1.3(b);
Lemma~2.4 and the endpoint formulas keep the affine rescaling and
the multipliers uniformly bounded after \(\mathcal U\) is fixed.
\end{proof}


\section{Local amplification with the external-node weight}
\label{sec:weighted}

Consider a sequence of finite node sets \(Y\subset[-1,1]\),
with \(N=\#Y\), and fix
\[
 I=(x_0-h,x_0+h)\Subset(-1,1),
 \qquad s=\#(Y\cap I).
\]
In the coordinate \(u=(x-x_0)/h\), define
\begin{equation}\label{eq:actual-field}
 W(u)=\prod_{y\in Y\setminus I}(x_0+hu-y),
 \qquad
 Q_s(u)=-\frac1s\log\left|\frac{W(u)}{W(0)}\right|.
\end{equation}
The function \(Q_s\) is convex on \((-1,1)\).

\begin{proposition}\label{prop:weighted}
For each \(H>1\), there are \(\epsilon_H,c_H,\delta_H>0\),
\(b_H\in(0,1)\), and a bounded complex neighborhood \(\Omega_H\)
of \([-b_H,b_H]\), with
\(\overline\Omega_H\subset\{z:|\operatorname{Re}z|<1\}\),
with the following property. Let \(\mathcal U_H\) be the ball
\(\norm{F-Q}_{\Omega_H}<\delta_H\) of bounded holomorphic
functions on \(\Omega_H\) that are real on \(\Omega_H\cap\R\).
For any fixed interval \(I\) and sequence of node sets \(Y\) as
above, suppose that \(s\to\infty\), that the normalized analytic
continuation of \(Q_s\) belongs to \(\mathcal U_H\), and that
\begin{equation}\label{eq:central-count}
 \#\{y\in Y:(y-x_0)/h\in A\}
 \le \frac{s}{4}+\epsilon_H s.
\end{equation}
If \(d\ge0\) are integers with \(d=o(s)\), then, for all sufficiently
large \(s\), there are
data \(v_y\in[-1,1]\), \(y\in Y\), such that every
\(p\in\Pi_{N+d}\) interpolating these data satisfies
\begin{equation}\label{eq:local-amplification}
 |\{x\in I:|p(x)|>H\}|\ge c_H|I|.
\end{equation}
\end{proposition}

\begin{proof}
We fix the auxiliary parameters in the following order, all before
letting \(s\to\infty\). First choose \(\rho\) so large that
\[
 c'\log\rho>10H,
\]
where \(c'\) is the constant in Lemma~\ref{lem:bernstein}, and set
\(\epsilon_H=1/(16\rho)\), \(R_0=\rho-\sqrt\rho\), and
\(\gamma=R_0^{-1/9}\). Choose \(t<1\)
so close to one that
\[
 \inf_{y\in A}2\pi\rho_t(y)>\pi-2\gamma.
\]
Fix \(K=3\), and then choose \(\kappa>0\) so small that
\(t+\kappa<1\) and the Jackson factor \eqref{eq:jackson} below is
larger than \(9/10\) whenever \(y\in A\) and
\(|c_s(u-y)|\le\rho+1\), for all sufficiently large \(s\). Choose
\(b\) with \(a_{t+\kappa}<b<1\), choose
\(A\Subset J_0\Subset J_1\Subset(-b,b)\) inside the common bulk
region, and shrink the analytic neighborhood of \(Q\) so that the
kernel estimates of Section~\ref{sec:kernel} are uniform there, the
supports of the equilibrium measures of masses near
\([t,t+\kappa]\) stay a fixed positive distance from \(\pm b\), and
\begin{equation}\label{eq:bandwidth-margin}
 \inf_{y\in A}\pi\beta_{s,y}>\pi-3\gamma.
\end{equation}
These choices give \(b_H=b\), \(\Omega_H\), \(\delta_H\), and
\(\mathcal U_H\) in the statement.
Set \(\eta=1/10\). By the uniform integral-tail bound in
Lemma~\ref{lem:bernstein}, choose \(T\) so that the truncation error
in \eqref{eq:kernel-projection} is less than \(\eta/2\). Next choose
\(R>\max\{\rho+1,2T\}\). Finally choose \(D>2R\) so that
\begin{equation}\label{eq:tail-choice}
 2C_\rho\sum_{j=0}^{\infty}
 \frac{1}{1+R+jD}\,
 \frac{1}{(1+\kappa(R+jD))^{2K}}<\frac1{10}.
\end{equation}
Here \(C_\rho\) is the constant in \eqref{eq:weighted-tail}; first
increasing \(R\) and then \(D\) makes \eqref{eq:tail-choice} hold.
All these parameters depend only on \(H\).

Partition the interval \(c_sA\) into boxes of length \(2\rho\),
discarding incomplete boxes at the ends.
There are \(s/(8\rho)-O(1)\) boxes.
Call a box admissible if it contains at most \(2\rho+1\) scaled nodes.
If \(B\) is the number of admissible boxes, then
\[
 \left(\frac{s}{8\rho}-O(1)-B\right)(2\rho+2)
 \le \frac{s}{4}+\epsilon_H s.
\]
The choice of \(\epsilon_H\) implies \(B\ge c_\rho s\).

For an admissible box centered at \(c_sy\), apply
Lemma~\ref{lem:bernstein} to its translated nodes, obtaining \(g\).
For a fixed truncation parameter \(T\), define
\begin{equation}\label{eq:kernel-projection}
 G_{s,y}(u)=\int_{-T}^{T}
 g(v)\frac{K_s(u,y+v/c_s)}{c_s}\,dv.
\end{equation}
This is \(e^{-sQ_s(u)}\) times a polynomial of degree at most \(m-1\).
By \eqref{eq:bandwidth-margin} and the Paley--Wiener theorem,
\[
 \int_{\R}g(v)
 \frac{\sin\bigl(\pi\beta_{s,y}(\xi-v)\bigr)}{\pi(\xi-v)}\,dv=g(\xi).
\]
With \(\eta,T,R\) fixed above, Equation~\eqref{eq:sine-limit}
implies, for all sufficiently large \(s\),
\begin{equation}\label{eq:projection-approximation}
 \sup_{|\xi|\le R}|G_{s,y}(y+\xi/c_s)-g(\xi)|<\eta.
\end{equation}
The uniform \(L^1\) bound makes this assertion uniform even when
\(g\) varies with the node configuration.

To obtain summable tails, multiply by an algebraic Jackson factor.
Write \(u=\cos\theta\), \(y=\cos\alpha\), and, with \(K=3\) and
\(\kappa\) fixed above, set \(k=\lfloor\kappa s/(2K)\rfloor\).
Define
\begin{align}
 D_k(v)&=\frac{\sin((k+1/2)v)}{(2k+1)\sin(v/2)},\notag\\
 L_{s,y}(u)&=D_k(\theta-\alpha)^{2K}
             +D_k(\theta+\alpha)^{2K}. \label{eq:jackson}
\end{align}
This expression is even in \(\theta\), so it is an algebraic
polynomial in \(u\) of degree at most \(2Kk\).
Since \(y\in A\), it is bounded on \([-1,1]\) by \(1+o(1)\).
By the choice of \(\kappa\), it is greater than \(9/10\) throughout
the scaled neighborhood containing the peak interval of \(g\).
Moreover,
\begin{equation}\label{eq:jackson-tail}
 |L_{s,y}(u)|\le
 C(1+\kappa s|u-y|)^{-2K}+Cs^{-2K}.
\end{equation}

Put \(F_{s,y}=G_{s,y}L_{s,y}\).
For \(u\in J_1\) with \(|c_s(u-y)|>2T\),
\eqref{eq:bulk-kernel} gives
\begin{equation}\label{eq:weighted-tail}
 \begin{split}
 |F_{s,y}(u)|
 &\le C_\rho(1+|c_s(u-y)|)^{-1}
       (1+\kappa|c_s(u-y)|)^{-2K}\\
 &\hspace{2em}+O_\rho(s^{-2K}).
 \end{split}
\end{equation}
For \(u\in[-b,b]\setminus J_0\),
\eqref{eq:edge-kernel} and \eqref{eq:jackson-tail} give
\begin{equation}\label{eq:weighted-edge}
 |F_{s,y}(u)|=O_\rho(s^{-2K-5/6}).
\end{equation}

To treat \(b<|u|<1\), write
\[
 F_{s,y}(u)=e^{-sQ_s(u)}H_{s,y}(u),
 \qquad
 \deg H_{s,y}\le D_s:=m-1+2Kk<s.
\]
Let \(\sigma\) be the equilibrium measure of mass \(D_s/s\)
for the field \(Q_s\) on \([-b,b]\).
Its support remains in a fixed compact subinterval of \((-b,b)\).
This follows from the continuity of the support endpoints in
\cite[Lemma 2.4]{KSSV}, after shrinking the external-field
neighborhood to cover masses near both \(t\) and \(t+\kappa\).
Let \(F\) be the constant value of \(Q_s-V_\sigma\) on
\(\supp\sigma\). Since \(\deg H_{s,y}\le D_s=s\sigma(\R)\), the
maximum principle applied to \(\log|H_{s,y}|-sV_\sigma\) on the
complement of the support, including infinity, yields
\begin{equation}\label{eq:weighted-maximum}
 |e^{-sQ_s(u)}H_{s,y}(u)|
 \le \norm{e^{-sQ_s}H_{s,y}}_{\supp\sigma}
 e^{-s(Q_s(u)-V_\sigma(u)-F)}.
\end{equation}
For \(u\in\supp\sigma\cap J_0\), the points
\(y+v/c_s\), \(|v|\le T\), lie in \(J_1\) for all large \(s\).
Hence \eqref{eq:bulk-kernel}, the uniform \(L^1\)-bound on \(g\),
and \(\norm{L_{s,y}}_\infty\le1+o(1)\) give
\[
 |F_{s,y}(u)|\le C_\rho.
\]
For \(u\in\supp\sigma\setminus J_0\), the same conclusion follows
from \eqref{eq:edge-kernel}, since \(y+v/c_s\) stays within
\(O(s^{-1})\) of \(A\). Consequently
\begin{equation}\label{eq:support-bound}
 \norm{e^{-sQ_s}H_{s,y}}_{\supp\sigma}
 =\norm{F_{s,y}}_{\supp\sigma}\le C_\rho
\end{equation}
uniformly in the admissible box and in the allowed external field.

Put \(D_s^*(u)=Q_s(u)-V_\sigma(u)-F\). The function \(Q_s\) is
convex on \((-1,1)\), while \(V_\sigma''<0\) off \(\supp\sigma\).
After shrinking the analytic neighborhood of \(Q\), there is
\(c_0>0\) such that \(Q_s''\ge c_0\) on \([-b,b]\), and hence
\[
 (D_s^*)''=Q_s''-V_\sigma''\ge c_0
 \qquad\text{off }\supp\sigma\text{ in }[-b,b].
\]
Let \(\alpha_s\) and \(\beta_s\) be the leftmost and rightmost points
of \(\supp\sigma\). The support remains a fixed positive distance
from \(\pm b\). Since \(D_s^*\ge0\) on \([-b,b]\) and \(D_s^*=0\) on the support,
its outward one-sided derivative at \(\alpha_s\) and \(\beta_s\) is
nonnegative. Integrating the preceding inequality from the support
to \(\pm b\) gives, for some \(c>0\) independent of \(s\),
\[
 D_s^*(-b)\ge c,\qquad D_s^*(b)\ge c.
\]
Outside \([-b,b]\), strict convexity on the two complementary
intervals shows that \(D_s^*\) continues to increase toward the
endpoints. Therefore \eqref{eq:weighted-maximum} and
\eqref{eq:support-bound} give
\begin{equation}\label{eq:weighted-exterior}
 |F_{s,y}(u)|=O_\rho(e^{-cs})
 \qquad(b<|u|<1).
\end{equation}

Retain admissible boxes whose centers in the \(c_su\) coordinate
are separated by at least the fixed distance \(D>2R\). By
\eqref{eq:tail-choice}, the sum of the tails in
\eqref{eq:weighted-tail} from all centers outside the \(R\)-window is
less than \(1/10\).
One residue class of the box indices retains at least
\(c_{\rho,D}s\) admissible boxes.
At any point there is at most one center within scaled distance \(R\).
In this window use \eqref{eq:projection-approximation}.
At nodes of the same box use \(|g|\le1\); outside the box use
\eqref{eq:bernstein-tail}. The remaining terms are controlled by
the absolutely summable tails
\eqref{eq:weighted-tail}--\eqref{eq:weighted-exterior}.
The \(O_\rho(s^{-2K})\) remainders remain \(o(1)\) after summing
over all retained boxes. Thus
\begin{equation}\label{eq:sum-node-bound}
 \max_{z\in Y\cap I}\sum_j
 \left|F_{s,y_j}\left(\frac{z-x_0}{h}\right)\right|\le2.
\end{equation}

On \((-1,1)\), the ratio \(W(u)/W(0)\) is positive, and hence
\(e^{-sQ_s(u)}=W(u)/W(0)\).
Each function
\[
 B_j(x)=F_{s,y_j}\left(\frac{x-x_0}{h}\right)
\]
therefore extends to a polynomial in the original coordinate \(x\),
vanishes at \(Y\setminus I\), and has degree at most
\begin{equation}\label{eq:degree-budget}
 N-s+m-1+2Kk\le N-1.
\end{equation}
On its peak interval, \eqref{eq:projection-approximation} and the
choice of \(\kappa\) give \(F_{s,y_j}\ge\frac9{10}(10H-\eta)\).
The other retained centers contribute at most \(1/10+o(1)\) there.
Thus the \(B_j\) satisfy the peak dominance and node bounds in
Lemma~\ref{lem:signs}, on node-free intervals of length at least
\(ch/s\) in the original coordinate. There are at least
\(c_{\rho,D}s\) such intervals. Since \(d=o(s)\), that lemma
supplies data for which every interpolant satisfies \(|p|>H\) on
a fixed fraction of them, proving \eqref{eq:local-amplification}.
\end{proof}

\section{Intervals with small node density}\label{sec:sparse}

\begin{proposition}[Low-density case]\label{prop:sparse}
For each \(H>1\), there is \(c_H>0\) with the following property.
Fix \(I\Subset(-1,1)\), and let \(I^\theta=\arccos I\subset(0,\pi)\).
Suppose that
\[
 \frac{\max_{\theta\in\overline{I^\theta}}\sin\theta}
      {\min_{\theta\in\overline{I^\theta}}\sin\theta}\le2.
\]
Let \(Y\subset[-1,1]\) be a sequence of sets with \(N=\#Y\to\infty\)
satisfying
\begin{equation}\label{eq:sparse-count}
 \#(Y\cap I)\le\frac{N|I^\theta|}{2\pi},
\end{equation}
and let \(d\ge0\) be integers with \(d=o(N)\).
For all sufficiently large \(N\), there are data \(v_y\in[-1,1]\)
on \(Y\) such that every \(p\in\Pi_{N+d}\) interpolating them satisfies
\begin{equation}\label{eq:sparse-amplification}
 |\{x\in I:|p(x)|>H\}|\ge c_H|I|.
\end{equation}
\end{proposition}

\begin{proof}
Write \(x=\cos\theta\) and \(v=N\theta/\pi\).
Choose \(\rho\) so that \(c'\log\rho>10H\), and then choose
\(D>4\rho\) so that
\begin{equation}\label{eq:sparse-tail-choice}
 2C\sum_{j=1}^{\infty}(jD-2\rho)^{-3}<\frac1{10}.
\end{equation}
Partition the corresponding \(v\)-interval into boxes of length
\(2\rho\), discarding the incomplete boxes at the two ends. Let \(M\)
be the number of complete boxes and \(B_{\rm bad}\) the number
containing at least \(2\rho+2\) nodes. The \(v\)-interval has length
\(L=N|I^\theta|/\pi\), so \eqref{eq:sparse-count} gives
\[
 B_{\rm bad}(2\rho+2)\le \frac L2,
 \qquad L<2\rho(M+2).
\]
Thus
\[
 B_{\rm bad}<\frac{\rho(M+2)}{2\rho+2}\le\frac M2+1.
\]
Hence at least \(M/2-1\) complete boxes contain at most \(2\rho+1\)
nodes; in particular, this is at least \(M/3\) once \(M\ge6\). Retain such boxes with a fixed separation \(D\),
and apply Lemma~\ref{lem:bernstein} at their centers \(a_j\).
For the resulting functions \(g_j\), define
\begin{equation}\label{eq:periodization}
 \Phi_{N,j}(\theta)=
 \sum_{\ell\in\Z}
 \left[
 g_j\left(\frac{N\theta}{\pi}-a_j+2N\ell\right)
 +g_j\left(-\frac{N\theta}{\pi}-a_j+2N\ell\right)
 \right].
\end{equation}
The series converges uniformly by \eqref{eq:bernstein-tail}.
Its Fourier coefficients vanish at all frequencies \(k\) for which
\[
 \frac{\pi|k|}{N}>\pi-4\gamma,
\]
by the Paley--Wiener theorem. It is therefore an even real
trigonometric polynomial of degree less than \(N\), and hence an
algebraic polynomial of degree less than \(N\) in \(x=\cos\theta\).

Because \(I^\theta\) stays away from \(0\) and \(\pi\), the reflected
centers stay a fixed angular distance from \([0,\pi]\), modulo
\(2\pi\). Their contribution, and that of every nonzero period of
the unreflected term, is \(O(N^{-3})\) uniformly on \([0,\pi]\), by
\eqref{eq:bernstein-tail}. Summing these errors over \(O(N)\)
retained boxes gives \(o(1)\).

By \eqref{eq:bernstein-tail} and
\eqref{eq:sparse-tail-choice}, the sum of the principal tails from all
other retained boxes is less than \(1/10\). For each
retained box let
\[
 B_{N,j}(x)=\Phi_{N,j}(\arccos x).
\]
If a node \(z=\cos\theta\) lies in the \(j\)-th retained box, the
principal term of \(B_{N,j}(z)\) has absolute value at most one by
\eqref{eq:bernstein-norm}; if it lies outside that box, the same term
is controlled by \eqref{eq:bernstein-tail}. The reflected and
nonzero-period terms contribute only \(o(1)\) after summation.
Therefore, for all sufficiently large \(N\),
\begin{equation}\label{eq:sparse-node-bound}
 \max_{z\in Y}\sum_j|B_{N,j}(z)|\le2.
\end{equation}
Each retained box contains a node-free interval of a fixed positive
length in the \(v\)-coordinate where its principal function is at
least \(c'\log\rho>10H\). The sum of all other terms is at most
\(1/10+o(1)\), so these intervals satisfy the peak dominance
condition of Lemma~\ref{lem:signs}.
There are at least \(c_{\rho,D}N|I^\theta|\) such intervals, each
of angular length at least \(c/N\). Since \(d=o(N)\), that lemma
supplies bounded data forcing every interpolant to satisfy
\(|p|>H\) on a fixed fraction of them. The assumed ratio bound for
\(\sin\theta\) converts their total angular length into at least
\(c_H|I|\) in the original coordinate, proving
\eqref{eq:sparse-amplification}.
\end{proof}

\section{Local intervals for an arbitrary node array}
\label{sec:arbitrary}

Fix \(H>1\). Throughout this section, \(S\subset[-1,1]\) denotes
an arbitrary fixed finite set, and \(Y_n=X_n\setminus S\).
Pass to a subsequence, with indices still denoted by \(n\), such that
\begin{equation}\label{eq:weak-limit}
 \nu_n=\frac1n\sum_{x\in X_n}\delta_x
 \ \Longrightarrow\ \nu.
\end{equation}
Only this subsequence will be used in the construction.
Removing the nodes belonging to any fixed finite set does not
alter the weak limit.
Put
\[
 V=V_\nu,\qquad
 m_0=\inf_{\supp\nu}V.
\]
The function \(V\) belongs to \(L^1[-1,1]\), and \(V\ge m_0\)
throughout \([-1,1]\).
To see the latter assertion off the support, note that \(V\) is
concave on each bounded complementary interval and has the
corresponding endpoint limits. On the two exterior intervals it
is monotone toward the support.

\subsection{Points above the minimum potential}
\label{subsec:high-potential}

Put \(P_n=P_{Y_n}\).
For \(z\in\supp\nu\), choose \(z_n\in Y_n\) with \(z_n\to z\).
Truncation of the logarithmic kernel gives
\begin{equation}\label{eq:potential-limits}
 \limsup_{n\to\infty}\frac1n\log|P_n'(z_n)|\le V(z),
 \qquad
 \frac1n\log|P_n|\longrightarrow V
 \quad\text{in }L^1[-1,1].
\end{equation}
For the first assertion, truncate \(\log|x-y|\) from below at
\(-A\). The contribution of the omitted diagonal term is \(A/n\);
one first takes the weak limit and then lets \(A\to\infty\).
For the second, use the continuity of
\[
 y\longmapsto\bigl(x\longmapsto\log|x-y|\bigr)
\]
as a map from \([-1,1]\) into \(L^1[-1,1]\).

Suppose first that \(V(z)>-\infty\), and let
\[
 A_{z,\eta}=\{x:V(x)>V(z)+3\eta\},\qquad \eta>0.
\]
Outside a subset of \(A_{z,\eta}\) whose measure tends to zero,
\eqref{eq:potential-limits} implies
\begin{equation}\label{eq:cardinal-growth}
 |\ell_n(x)|\ge e^{\eta n},
 \qquad
 \ell_n(x)=\frac{P_n(x)}{(x-z_n)P_n'(z_n)}.
\end{equation}
Specify the value one at \(z_n\), and zero at all other nodes of
\(Y_n\). Every polynomial of degree at most \(n+r_n\) satisfying
these data can be written as
\begin{equation}\label{eq:cardinal-correction}
 p=\ell_n T,\qquad T(z_n)=1,\qquad
 \deg T\le r_n+\#S+1=o(n).
\end{equation}
If \(|p|\le H\) on a subset of \(A_{z,\eta}\) of fixed positive
measure, \eqref{eq:cardinal-growth} and Remez's inequality
\cite[p.~93, (3)--(4)]{Remez36} yield
\[
 |T(z_n)|\le He^{-\eta n}C^{\,r_n+\#S+1}<1
\]
for all sufficiently large \(n\), a contradiction.
Here \(C\) depends only on the fixed measure lower bound.
Thus the measure of the low set in \(A_{z,\eta}\) tends to zero
uniformly over all the polynomials in \eqref{eq:cardinal-correction}.
At every density point of \(A_{z,\eta}\), each sufficiently small
fixed interval therefore admits data forcing large values on at
least half its length.

Choose countably many \(z\) for which \(V(z)\) decreases to \(m_0\),
and take positive rational \(\eta\).
The sets \(A_{z,\eta}\) cover \(\{V>m_0\}\) when the chosen source
potentials are finite.
If some \(V(z)=-\infty\), the first assertion of
\eqref{eq:potential-limits} instead gives
\[
 \frac1n\log|P_n'(z_n)|\longrightarrow-\infty.
\]
On each set \(\{V\ge-A\}\), the same argument gives
\eqref{eq:cardinal-growth} for any fixed positive exponential rate,
outside a set of vanishing measure. Letting \(A\to\infty\) covers
almost all of \([-1,1]\).
In particular, if \(m_0=-\infty\), this subsection supplies local
intervals at almost every point.

\subsection{The minimum-potential set}
\label{subsec:minimum}

Suppose that \(m_0>-\infty\), and let
\[
 E=\{x:V(x)=m_0\}.
\]
By the differentiation theorem for measures, \(\nu\) has a finite
Lebesgue density at almost every point.
At points where this density is zero, all sufficiently small
intervals satisfy the strict limiting form of
\eqref{eq:sparse-count}. Fixing such an interval before taking
\(n\to\infty\), and choosing endpoints that are not atoms of \(\nu\),
allows Proposition~\ref{prop:sparse} to be applied.

We isolate the local field arising at the remaining points.

\begin{lemma}[Local outer field]\label{lem:local-outer-field}
Let \(x_0\) be a density point of \(E\) and a Lebesgue point of the
absolutely continuous part of \(\nu\), with finite positive density
\(w\), and suppose that the singular part of \(\nu\) has density zero
at \(x_0\). Put
\[
 I_h=(x_0-h,x_0+h),\qquad a_h=\nu(I_h),
\]
and let \(\mu_h\) be the image of \(\nu|_{I_h}/a_h\) under
\(x\mapsto(x-x_0)/h\). Then, as \(h\downarrow0\),
\[
 a_h\sim2wh,\qquad \mu_h\Longrightarrow \frac{du}{2},
\]
and
\begin{equation}\label{eq:inner-potential}
 V_{\mu_h}\longrightarrow Q-1
 \quad\text{in }L^1(-1,1).
\end{equation}
Define the normalized outer potential
\[
 A_h(u)=-\frac1{a_h}
 \int_{\R\setminus I_h}\log|x_0+hu-y|\,d\nu(y).
\]
Then
\begin{equation}\label{eq:real-field-limit}
 A_h(u)-A_h(0)\longrightarrow Q(u)
 \quad\text{locally uniformly on }(-1,1).
\end{equation}
Moreover, the functions \(A_h-A_h(0)\) admit analytic continuations
\(\widetilde A_h\) to the strip
\(\Omega_0=\{z\in\C:|\operatorname{Re}z|<1\}\), and
\[
 \widetilde A_h\longrightarrow Q
\]
locally uniformly there.
\end{lemma}

\begin{proof}
The differentiation theorem gives \(a_h\sim2wh\) and
\(\mu_h\Rightarrow du/2\). The continuity of
\[
 y\longmapsto\bigl(x\longmapsto\log|x-y|\bigr)
\]
from \([-1,1]\) into \(L^1[-1,1]\) gives
\eqref{eq:inner-potential}, since the logarithmic potential of
\(du/2\) is \(Q-1\).

For \(u\) corresponding to \(E\cap I_h\), the equality \(V=m_0\)
gives
\begin{equation}\label{eq:outer-identity}
 A_h(u)+\frac{m_0}{a_h}-\log h=V_{\mu_h}(u).
\end{equation}
The relative Lebesgue measure of this set in \((-1,1)\) tends
to one, and the left side of \eqref{eq:outer-identity} is convex.
Thus \eqref{eq:inner-potential} and \eqref{eq:outer-identity} imply
convergence in measure to \(Q-1\). Finite convex functions that
converge in measure to a continuous finite convex function converge
uniformly on compact subintervals: convergence points on either side
of a compact interval bound the secant slopes, giving equicontinuity.
This proves \eqref{eq:real-field-limit}.

Writing \(\xi=(y-x_0)/h\), one has \(|\xi|\ge1\) for
\(y\notin I_h\). Define \(\widetilde A_h\) on \(\Omega_0\) by
\(\widetilde A_h(0)=0\) and
\[
 \widetilde A_h'(z)
 =-\frac{h}{a_h}\int_{\R\setminus I_h}
   \frac{d\nu(y)}{x_0+hz-y}.
\]
For real \(u\in(-1,1)\), integration along the real segment gives
\(\widetilde A_h(u)=A_h(u)-A_h(0)\), and
\[
 \widetilde A_h''(z)
 =\frac1{a_h}\int_{\R\setminus I_h}\frac{d\nu(y)}{(z-\xi)^2}.
\]
Fix a compact subset \(K\) of \(\Omega_0\), and choose a compact
interval \(K_1\Subset(-1,1)\) whose interior contains
\(\operatorname{Re}K\cup\{0\}\).
For sufficiently small fixed \(\ell>0\), the kernels
\((u-\xi)^{-2}\) are uniformly comparable on
\([u-\ell,u+\ell]\), uniformly for \(u\in K_1\). Hence
\[
 A_h''(u)
 \le C_\ell\int_{u-\ell}^{u+\ell}A_h''(v)\,dv
 =C_\ell\bigl(A_h'(u+\ell)-A_h'(u-\ell)\bigr).
\]
The local uniform convergence in \eqref{eq:real-field-limit} and
convexity give uniform interior bounds for \(A_h'\), hence for
\(A_h''\) on \(K_1\). Since
\[
 |z-\xi|^{-2}\le|\operatorname{Re}z-\xi|^{-2},
\]
\(\widetilde A_h''\) is uniformly bounded for
\(\operatorname{Re}z\in K_1\). Integrating from zero and using the
interior first-derivative bound gives local boundedness on \(\Omega_0\),
hence a normal family. Every subsequential limit agrees
with \(Q\) on the real interval by \eqref{eq:real-field-limit}, and
therefore equals \(Q\) by the identity theorem.
\end{proof}

\begin{lemma}[Discrete outer-field approximation]\label{lem:discrete-field}
Fix \(h>0\) such that the endpoints of \(I_h\) and the preimages of
the endpoints of \(A\) are not atoms of \(\nu\), and put
\(s_n=\#(Y_n\cap I_h)\).
Then
\begin{equation}\label{eq:discrete-counts}
 \frac{s_n}{n}\longrightarrow a_h,
 \qquad
 \frac{\#\{y\in Y_n:(y-x_0)/h\in A\}}{s_n}
 \longrightarrow\mu_h(A).
\end{equation}
Let \(\mathcal Q_n\) be the normalized analytic continuation of the
discrete field \eqref{eq:actual-field}, with \(\mathcal Q_n(0)=0\).
On every fixed complex neighborhood compactly contained in the domain
of \(\widetilde A_h\),
\[
 \mathcal Q_n\longrightarrow\widetilde A_h
\]
locally uniformly. If \(N_n=\#Y_n\), then the degree bound \(n+r_n\)
is \(N_n+d_n\), where
\begin{equation}\label{eq:discrete-excess}
 d_n=r_n+\#(X_n\cap S)=o(s_n).
\end{equation}
\end{lemma}

\begin{proof}
The two limits in \eqref{eq:discrete-counts} follow from weak
convergence and the choice of the interval endpoints. Moreover,
\[
 \mathcal Q_n'(z)
 =-\frac{h}{s_n}\sum_{y\in Y_n\setminus I_h}
   \frac1{x_0+hz-y}.
\]
On a fixed complex neighborhood as above, the summands are uniformly
bounded and equicontinuous as functions of the outer node. Since
\(n/s_n\to1/a_h\), weak convergence of the empirical measures gives
\(\mathcal Q_n'\to\widetilde A_h'\) locally uniformly. Integrating
from zero gives the asserted convergence of the fields. Finally,
\(r_n=o(n)\) and \(s_n/n\to a_h>0\) give
\eqref{eq:discrete-excess}.
\end{proof}

Now let \(x_0\) satisfy the hypotheses of Lemma~\ref{lem:local-outer-field};
these conditions hold at almost every point of \(E\) where the density
of \(\nu\) is positive. Since \(\overline\Omega_H\Subset\Omega_0\),
Lemma~\ref{lem:local-outer-field} gives
\(\norm{\widetilde A_h-Q}_{\Omega_H}\to0\). Choose
\(h\) so small that this norm is less than \(\delta_H/2\), for the
neighborhood \(\mathcal U_H\) in Proposition~\ref{prop:weighted}, and
\[
 \mu_h(A)<\frac14+\frac{\epsilon_H}{2}.
\]
This is possible because \(\mu_h\Rightarrow du/2\) and
\((du/2)(A)=1/4\). We may also choose \(h\) so that the interval
endpoints required in Lemma~\ref{lem:discrete-field} are not atoms of
\(\nu\). Lemma~\ref{lem:discrete-field} gives
\(\norm{\mathcal Q_n-\widetilde A_h}_{\Omega_H}<\delta_H/2\) for all
sufficiently large \(n\), so \(\mathcal Q_n\in\mathcal U_H\).
It also gives \eqref{eq:central-count} and an excess degree
\(d_n=o(s_n)\). Therefore
Proposition~\ref{prop:weighted} applies.

\begin{proposition}[Local forcing for arbitrary arrays]\label{prop:local-forcing}
For each \(H>1\), there is \(c_H>0\) such that every array
\((X_n)\) and sequence \((r_n)\) as in Theorem~\ref{thm:main}
admit a subsequence with the following property.
For almost every \(x\in(-1,1)\), there are arbitrarily small
intervals \(I\Subset(-1,1)\) containing \(x\) such that, for every
fixed finite \(S\subset[-1,1]\), all sufficiently late rows of the
subsequence admit data bounded by one on \(X_n\setminus S\) forcing
\[
 |\{y\in I:|p_n(y)|>H\}|\ge c_H|I|
\]
for every interpolant of degree at most \(n+r_n\).
\end{proposition}

\begin{proof}
Use the weakly convergent subsequence in \eqref{eq:weak-limit}.
Section~\ref{subsec:high-potential} gives the assertion on
\(\{V>m_0\}\), with fraction \(1/2\). On \(\{V=m_0\}\),
Section~\ref{subsec:minimum} applies Proposition~\ref{prop:sparse}
where the density of \(\nu\) is zero, and
Proposition~\ref{prop:weighted} where it is positive and finite.
These cases exhaust almost every point. Take the minimum of their
three constants. The intervals are selected from \(\nu\) and \(H\);
for every fixed \(S\), the same limits apply to \(X_n\setminus S\),
with excess degree \(r_n+\#(X_n\cap S)\).
\end{proof}

\section{A finite construction for common low sets}
\label{sec:finite}

\begin{proposition}\label{prop:finite}
Fix \(H>1\), \(\delta>0\), and \(N_0\ge1\).
There are finitely many indices \(n_1,\ldots,n_k\ge N_0\) and
a function \(g\in C([-1,1];\R)\), with \(\norm{g}_\infty\le1\),
such that every choice of polynomials
\[
 p_j\in\Pi_{n_j+r_{n_j}},\qquad p_j=g\text{ on }X_{n_j},
\]
satisfies
\begin{equation}\label{eq:finite-low-set}
 \left|\bigcap_{j=1}^k\{x:|p_j(x)|\le H\}\right|<\delta.
\end{equation}
\end{proposition}

\begin{proof}
For any interval \(J\subset(-1,1)\), the intervals furnished by
Proposition~\ref{prop:local-forcing} form a Vitali cover of almost all of \(J\).
Choose finitely many disjoint such intervals whose total length is
at least \(|J|/2\).
For each interval, choose a sufficiently late new row of the fixed
subsequence, assigning values only at nodes whose values have not
yet been specified, applying Proposition~\ref{prop:local-forcing}
with \(S\) equal to the finite set of previously assigned nodes.
For all choices of the resulting interpolating polynomials, the
union of their sets of values exceeding \(H\) occupies at least
\begin{equation}\label{eq:interval-union}
 c_*|J|,\qquad c_*=\frac{c_H}{2}>0,
\end{equation}
inside \(J\).

Assume \(0<\delta\le2\), since larger \(\delta\) is immediate.
Suppose that finitely many rows have already been chosen, with
degree bounds \(d_1,\ldots,d_k\).
For every admissible choice of their polynomials, the common low set
\[
 L=\bigcap_{j=1}^k\{x:|p_j(x)|\le H\}
\]
has at most \(2+2\sum_jd_j\) boundary points in \([-1,1]\).
A constant polynomial equal to \(H\) or \(-H\) introduces no
additional boundary points.
Choose a fixed partition of \([-1,1]\) into sufficiently short
intervals that the cells meeting these boundary points have total
length less than \(\delta/4\).
The same partition works for every admissible choice of the
polynomials, since the boundary bound depends only on their degrees.

In each cell, carry out the construction giving
\eqref{eq:interval-union}, using new rows and assigning only new
node values.
If \(|L|\ge\delta\), the cells contained entirely in \(L\) have
total length at least
\[
 |L|-\delta/4\ge3|L|/4.
\]
For the new common low set this gives
\begin{equation}\label{eq:low-contraction}
 |L_{\mathrm{new}}|
 \le (1-3c_*/4)|L|
 \qquad\text{whenever }|L|\ge\delta.
\end{equation}
Starting from \(L=[-1,1]\), perform a finite number \(q\) of batches,
where
\[
 2(1-3c_*/4)^q<\delta.
\]
At each batch the new partition is chosen using the degree bounds of
the rows already selected.
Equation~\eqref{eq:low-contraction} shows that the final common
low set has measure less than \(\delta\), uniformly over all
admissible polynomial choices.
All rows may be chosen beyond \(N_0\).

The assigned values define a function on a finite union of nodes,
with absolute value at most one. Extend it by linear interpolation
between consecutive nodes and by constants to the endpoints of
\([-1,1]\). The resulting continuous function \(g\) has
\(\norm{g}_\infty\le1\) and satisfies \eqref{eq:finite-low-set}.
\end{proof}

\section{Proof of Theorem~\ref{thm:main}}\label{sec:baire}

Work in the real Banach space \(C([-1,1];\R)\).
For \(H>0\) and positive integers \(N,k\), let \(\mathcal O_{H,N,k}\)
consist of the functions \(f\) for which there is a finite set
\(S\subset\{N,N+1,\ldots\}\) such that every choice
\[
 p_n\in\Pi_{n+r_n},\qquad p_n=f\text{ on }X_n\quad(n\in S),
\]
satisfies
\[
 \left|\bigcap_{n\in S}\{x:|p_n(x)|\le H\}\right|<\frac1k.
\]

\begin{lemma}\label{lem:stability}
If \(H>M>0\), then
\(\mathcal O_{H,N,k}\subset\operatorname{int}\mathcal O_{M,N,k}\).
\end{lemma}

\begin{proof}
Let \(f\in\mathcal O_{H,N,k}\), with finite set \(S\) as above.
For each \(n\in S\), let \(L_n\) be the Lagrange interpolation
operator on \(X_n\), and write \(\ell_{n,\xi}\) for its fundamental
polynomials. Put
\[
 C=1+\sum_{n\in S}\sum_{\xi\in X_n}\norm{\ell_{n,\xi}}_\infty.
\]
If \(\norm{g-f}_\infty<(H-M)/C\) and \(p_n\) interpolates \(g\),
then \(q_n=p_n-L_n(g-f)\) interpolates \(f\) with the same degree
bound. Moreover,
\[
 \norm{L_n(g-f)}_\infty<H-M,
 \qquad
 \{|p_n|\le M\}\subset\{|q_n|\le H\}.
\]
The defining bound for \(f\) therefore gives the defining bound for
\(g\) at height \(M\). Thus the ball of radius \((H-M)/C\) about
\(f\) lies in \(\mathcal O_{M,N,k}\).
\end{proof}

\begin{lemma}\label{lem:dense-open}
For every positive integer triple \(M,N,k\), the open set
\(\operatorname{int}\mathcal O_{M,N,k}\) is dense.
\end{lemma}

\begin{proof}
Let an open ball in \(C([-1,1];\R)\) be given.
By polynomial approximation, it contains a polynomial \(h\).
Choose \(\varepsilon>0\) so that \(h+\varepsilon g\) belongs to
the ball whenever \(\norm{g}_\infty\le1\).
Apply Proposition~\ref{prop:finite} with
\[
 H>\max\left\{1,\frac{M+1+\norm{h}_\infty}{\varepsilon}\right\},
 \qquad \delta=\frac1k,
\]
requiring all selected indices to be at least \(N\) and large enough
that \(n+r_n\ge\deg h\).
Let \(g\) be the resulting function and set \(f=h+\varepsilon g\).
For any admissible interpolant \(p_n\) of \(f\) on a selected row,
\[
 q_n=\frac{p_n-h}{\varepsilon}
\]
is an admissible interpolant of \(g\), and \(|p_n|\le M+1\)
implies \(|q_n|<H\). Hence \(f\in\mathcal O_{M+1,N,k}\).
Lemma~\ref{lem:stability} puts \(f\) in
\(\operatorname{int}\mathcal O_{M,N,k}\), as required.
\end{proof}

\begin{proof}[Proof of Theorem~\ref{thm:main}]
By Lemma~\ref{lem:dense-open} and the Baire category theorem, choose
\[
 f\in\bigcap_{M,N,k\ge1}\operatorname{int}\mathcal O_{M,N,k}.
\]
Fix any admissible interpolation sequence \((p_n)\).
For every \(M,N,k\), the common tail low set
\[
 E_{M,N}=\bigcap_{n\ge N}\{x:|p_n(x)|\le M\}
\]
is contained in the finite common low set supplied by
\(f\in\mathcal O_{M,N,k}\). Thus \(|E_{M,N}|<1/k\) for every
\(k\), and \(|E_{M,N}|=0\).
Outside the null set \(\bigcup_{M,N\ge1}E_{M,N}\), every tail
of \((|p_n(x)|)\) exceeds every positive integer threshold.
This proves \eqref{eq:divergence}.
\end{proof}

\section*{Code availability}

A partial Lean~4 formalization is available at
\url{https://github.com/FireflySentinel/erdos-1152}.
It covers the sign-change bound in Lemma~\ref{lem:signs}, the scalar
calculations of Lemma~\ref{lem:equilibrium}, the deduction from
\eqref{eq:cardinal-growth} using Remez's inequality, and the finite
construction and Baire category argument in
Sections~\ref{sec:finite}--\ref{sec:baire}.
Remez's inequality and the remaining analytic estimates are hypotheses.

\section*{Use of generative AI}

The author used GPT-5.6 and GPT-6 Astra in developing the arguments
involving polynomial corrections, local external fields, weighted kernels
and finite interpolation, and OpenAI Codex (GPT-6) for the Baire category
argument and the Lean formalization. The author checked the arguments and
their use of the cited results and is responsible for the content.

\begin{thebibliography}{9}

\bibitem{EV80}
P.~Erd\H{o}s and P.~V\'ertesi,
On the almost everywhere divergence of Lagrange interpolatory polynomials
for arbitrary system of nodes,
\emph{Acta Math. Acad. Sci. Hungar.} \textbf{36} (1980), 71--89.
\url{https://www.renyi.hu/~p_erdos/1981-03.pdf}.

\bibitem{EV81}
P.~Erd\H{o}s and P.~V\'ertesi,
Correction of some misprints in our paper ``On the almost everywhere
divergence of Lagrange interpolatory polynomials for arbitrary system of
nodes'',
\emph{Acta Math. Acad. Sci. Hungar.} \textbf{38} (1981), 263.

\bibitem{EKS89}
P.~Erd\H{o}s, A.~Kro\'o, and J.~Szabados,
On convergent interpolatory polynomials,
\emph{J. Approx. Theory} \textbf{58} (1989), 232--241.
\url{https://www.renyi.hu/~p_erdos/1989-16.pdf}.

\bibitem{Favorite99}
B.~Bollob\'as et al. (collectors),
\emph{Some of Paul's Favorite Problems},
booklet circulated at the conference \emph{Paul Erd\H{o}s and His Mathematics},
Budapest, July 1999, Problem~2.42.
\url{https://web.math.pmf.unizg.hr/~vjekovac/EP/Some_of_Pauls_favorite_problems.pdf}.

\bibitem{Erdos1152}
T.~F. Bloom,
\emph{Erd\H{o}s Problem \#1152},
\url{https://www.erdosproblems.com/1152}.
Accessed 4 September 2026.

\bibitem{SV90}
J.~Szabados and P.~V\'ertesi,
\emph{Interpolation of Functions},
World Scientific, Singapore, 1990.

\bibitem{Remez36}
E.~Remes,
Sur une propri\'et\'e extr\'emale des polynomes de Tchebychef,
\emph{Comm. Inst. Sci. Math. M\'ec. Univ. Kharkoff Soc. Math. Kharkoff}
(4) \textbf{13} (1936), no.~1, 93--95.
\url{https://history-of-approximation-theory.com/fpapers/remezppr.pdf}.

\bibitem{KSSV}
T.~Kriecherbauer, K.~Schubert, K.~Sch\"uler, and M.~Venker,
Global asymptotics for the Christoffel--Darboux kernel of random
matrix theory,
\emph{Markov Processes and Related Fields} \textbf{21} (2015),
639--694.
\href{https://arxiv.org/abs/1401.6772}{arXiv:1401.6772}.

\bibitem{OU}
A.~Olevskii and A.~Ulanovskii,
\href{https://doi.org/10.1134/S0081543818080151}
{On irregular sampling and interpolation in Bernstein spaces},
\emph{Proceedings of the Steklov Institute of Mathematics}
\textbf{303} (2018), 178--192.

\end{thebibliography}

\end{document}
